\clearpage
\section{Anexo 8.E: Controles técnicos preventivos}
Los 15 controles vigentes (CT-01--CT-15) establecen obligaciones de diseño, innovación y calidad de LafroX. CT-03 reúne identidad conectada y local conforme a S4/S9; CT-15 conserva stock, antes identificado como CT-16. Los 15 identificadores vigentes se mantienen; la ampliación de identidad no crea otra contingencia ni renumera stock. Cada ficha define una solución técnica, su responsable y la evidencia necesaria para verificarla. La habilitación requiere ejecutar las pruebas y aprobar la evidencia de cumplimiento. CT sin número designa la Contraparte Técnica.

CAL verifica la evidencia y JP controla el cierre. SRE asume el trabajo operativo de INN-03 (familia DAT) desde el mes 21 y dirige el despliegue por rutas de INN-05 (familia IMP) desde el mes 22. La Tabla 8.1 identifica los responsables por familia. No se promueve un flujo con defectos críticos o altos abiertos ni con pruebas habilitantes fallidas. Una función opcional que no supere sus pruebas permanece deshabilitada y se conserva el alcance obligatorio.

\begingroup
\renewcommand{\thetable}{8.E.1}
\begin{landscape}
\setlength{\textwidth}{\linewidth}
\renewcommand{\arraystretch}{1.08}
\begin{tablalafrox}{Controles técnicos: solución, responsable y verificación. Fuente: requisitos de arquitectura, seguridad, innovación y calidad de la solución.}{tab:cierre-ci}{G{0.08} F{0.25} Y F{0.34}}{\cab{Control} & \cab{Responsable} & \cab{Solución requerida} & \cab{Evidencia y condición de cierre}}
CT-01 & ARQ; DES implementa; CAL verifica. & Idempotencia persistente, consulta de estado antes de reintentar, conciliación y estado pendiente sin confirmar. & Simular timeout antes/después del cargo y reenvío de la misma clave: una sola imputación y conciliación verificable. Antes de H5/H10 en el alcance aplicable; CAL bloquea promoción del flujo fallido. \\
CT-02 & ARQ; SRE opera; CAL verifica. & Definir circuit breaker, caché de rutas autorizadas, límites y estado degradado ante caída del servicio de mapas elegido, sin dar nuevas rutas por calculadas. & Inyectar caída y recuperación del proveedor: continuidad autorizada de rutas disponibles, nuevas consultas pendientes y restitución sin pérdida. Antes de H5/H10 y tras cambio de proveedor; corregir el diseño si falla. \\
CT-03 & SEG responde; ARQ diseña; DES/SRE ejecutan; CAL verifica. & Autenticación y sesión conectadas y locales, mínimo privilegio, detección, revocación y auditoría. Manifiesto de bodega de 26 h; al aislar una hora después quedan 25 h. Verificar relevos a las 8/16 h, firma, sitio, dispositivo, rol, PIN y turno. Reinicio/reloj no extienden vigencia. Reparto de 14,5 h con autorización máxima de 16 h desde emisión. & AL-ID-01 y AL-OFF-01 en resiliencia 6.3.1/6.3.2 antes de H5/H10; seguridad sin hallazgos críticos/altos antes de G5. Exigir cero accesos indebidos, cola conservada y conciliación única. Distinguir baja local inmediata y baja central durante aislamiento; verificar el detalle de esta ficha y 8.D.3. Corregir y repetir antes de habilitar. \\
CT-04 & ARQ; SEG valida autenticidad; CAL verifica. & Validar origen, integridad, correlación y repetición antes de procesar EDI; poner mensajes inválidos en cuarentena sin producir pedidos o avisos. & Pruebas negativas con firma/origen alterados, mensaje repetido y correlación inválida: rechazo trazable sin efecto de negocio. Antes de H10 para EDI y regresión del alcance aplicable; corregir antes de promover. \\
CT-05 & SRE; SEG define política; CAL verifica. & Proteger almacenamiento y sesiones, aplicar enrolamiento y capacidades de bloqueo/gestión remota disponibles; revocar accesos tras pérdida y delimitar exposición. & Ensayo de terminal extraviado en línea y sin señal: datos protegidos, acceso no autorizado denegado y revocación/bloqueo al recuperar enlace. Antes de H5/H10 y tras cambio de política; no presumir borrado remoto sin conectividad. \\
CT-06 & SEG; SRE ejecuta; CAL verifica. & Aplicar límites, protección perimetral, aislamiento de portales y alertas; conservar capacidad para usuarios legítimos y continuidad de servicios críticos. & Inyectar ráfaga controlada en ambiente autorizado y medir bloqueo, latencia y disponibilidad contra umbrales de S4/T-13. Antes de H5/H10 y tras cambio perimetral; ningún ensayo ofensivo no autorizado en Producción. \\
CT-07 & SEG; SRE ejecuta; CAL verifica. & Restringir privilegios, separar funciones, preservar auditoría y respaldos y comprobar recuperación ante cambios no autorizados de una cuenta administrativa. & Ensayar cambio indebido, detección, contención y restauración con integridad demostrada; recuperar conforme RTO/RPO exigidos. Antes de H5/H10, restauraciones y pruebas de seguridad periódicas de S9; no cerrar con evidencia comprometida. \\
CT-08 & DAT dirige; IMP coordina piloto; CAL verifica. & Validar catálogo con Comercial/Calidad en 7--9, mostrar confianza y confirmar saldo en visita. Mantener alcance obligatorio; si no se valida, usar sólo aviso informativo por vida útil remanente. & Paquete 11--15: H5 en 11--12 y piloto real en 13--15. I-01A en mes 15: al menos 70\% de aciertos sobre al menos 100 avisos evaluables. Sin cumplir, no ampliar INN-01 a otras rutas; corregir y repetir. Registrar los resultados del piloto. \\
CT-09 & CAL; DES implementa; SRE aporta incidentes. & Comparar reproducción con incidente de terreno y limitar captura de batería/datos. Conservar diagnóstico convencional y regresión base si la función opcional no demuestra utilidad. & Conservar casos incidente--reproducción y medición de I-02A y consumo contra lo comprometido en S13, antes de ampliar captura; corregir captura/reproducción fallida y no acreditar el indicador sin evidencia. \\
CT-10 & SEG; DES implementa; CAL verifica. & Sanitizar por reglas fijas antes de sacar paquetes del terminal; minimizar datos y modelar amenazas. Suspender conjuntos afectados y producir muestra protegida ante filtración. & Pruebas con datos personales sembrados y campos no previstos: ningún dato personal sin sanitizar abandona el terminal; informe de revisión de amenazas y regresión. Bloquear la captura/exportación hasta aprobar, antes del primer uso y ante cada cambio de esquema. \\
CT-11 & DAT; Calidad del CLIENTE valida parámetros; CAL verifica. & Validar familias, mantener RNG-04 y no extender vencimiento por inferencia; presentar la estimación sólo en M9. Familia no validada conserva fecha impresa e historia visible. & Evidencia de I-03B por familia conforme a S13 antes de habilitar orden sugerido, con pruebas de no relajación de RNG-04. CAL bloquea familia no validada, DAT corrige y repite; la participación del CLIENTE no traslada responsabilidad de diseño. \\
CT-12 & DAT; DES implementa; CAL verifica. & Vincular lote--sensor mediante eventos de ubicación M2 y declarar tramos sin dato; impedir estimación cuando la historia esté incompleta. & Ensayar movimiento sin evento, sensor equivocado y sincronización tardía: historia incompleta visible y ninguna vida útil inferida sin fundamento. Antes de habilitar INN-03 y en regresión; corregir asociación antes de promover. \\
CT-13 & IMP; DAT/DES implementan; SEG revisa privacidad; CAL verifica. & Codiseñar hoja comprensible; aplicar umbral de agregación y supresión en cada emisión comparativa y revisión legal previa. Conservar versión de un bloque si no se entiende y bloquear bloque 4 si expone terceros. & Prueba I-05B con usuarios y pruebas de inferencia conforme S13 antes de habilitar cada bloque; informe de comprensión y supresión. Corregir y repetir antes de habilitar. \\
CT-14 & SRE; ARQ dimensiona; DES corrige; CAL verifica. & Comparar perfil de carga acordado y observado, concentración por ventana y mezcla de operaciones; demostrar capacidad efectiva por dependencia, incluidos ERP, colas y recursos compartidos, peak conocido y crecimiento comprometido. Ajustar el dimensionamiento ofertado sin degradar requisitos. & Pruebas de carga/estrés y p95 según S4/T-13, incluyendo peak y cuello de botella ERP, antes de H5/H10 y preparación de peak. Medir capacidad, colas, errores y recuperación; bloquear promoción si incumple. Conservar perfil, entradas, resultados y aprobación. Ante prueba fallida, LafroX corrige y repite antes de promover; no presume demanda adicional del CLIENTE. R-036 trata divergencias del perfil; R-031, coordinación del cierre; R-030, interfaces legadas. \\
CT-15 & Equipo DAT hasta mes 21; SRE desde mes 22; ARQ diseña; CAL verifica. & Una reserva central confirmada exige retención local durable de igual UUID, época y cantidad. Mantener una sola autoridad de escritura por ámbito; excluir escritor previo al promover, conciliar sitio/UUID y conservar capturas pendientes sin confirmar desde copia vencida. & AL-STOCK-01: última unidad concurrente, acuse perdido, reintento, cancelación y época obsoleta, sin doble reserva ni stock negativo; preventa real E1 (6.1.1), portal/EDI reales E2 (6.1.2). AL-DR-01: pérdida, promoción, retorno y conciliación (6.4.1/6.4.2), antes de H5/H10. CAL bloquea flujo inconsistente y exige corrección/regresión; simuladores no acreditan producto E2. \\
\end{tablalafrox}
\end{landscape}
\endgroup

\subsection*{Referencias de diseño y pruebas previstas}
\addcontentsline{toc}{subsection}{Referencias de diseño y pruebas previstas}

Las referencias siguientes localizan el sustento técnico de cada control en las versiones corregidas de S4, S9 y S13. Identifican diseño y verificación por ejecutar, no resultados aprobados. La ficha de 8.E.1 conserva la condición de cierre. Cuando sólo existe cobertura general, se declara la dependencia para que la rama responsable complete el protocolo específico y su trazabilidad antes de habilitar el flujo.

\begin{description}
\item[CT-01.] S4, «Módulo M7 Cobranza y rendición», en «Módulos funcionales y límites de contexto»: consulta del pago incierto por la misma clave, sin nuevo cobro ni autorización por silencio. T-13 de S9, 7.2, contempla la dependencia externa sin respuesta, con error o lenta; es cobertura general, no un caso POS específico acreditado. ARQ/DES deben precisar el contrato con la pasarela en S4 y CAL incorporar en S9 los casos de timeout, reenvío y conciliación de la ficha.
\item[CT-02.] S4, Anexo 4-D, «Módulo M4 Rutas», integración con GIS: API externa y caché de mapas por zona. La continuidad sin conexión que allí se describe pertenece al diseño previo y no se reincorpora como solución vigente. T-13, 7.2, aporta sólo el escenario general de dependencia externa; no acredita un protocolo de circuit breaker. ARQ/SRE deben completar en S4 los límites y estados degradados del diseño conectado y CAL su caída/recuperación específica en S9.
\item[CT-03.] S4, Anexo 4-M, AL-ID-01 y AL-OFF-01, y Anexo 4-R, controles 5.15/5.18, 8.5 y 8.15/8.16: identidad conectada y local, permisos, sesión, bajas y auditoría. S9, T-13, 7.5, asigna resiliencia 6.3.1/6.3.2 y define el expediente de AL-ID-01; T-17, 3.1, conserva versión y actas. La ficha incluye esa continuidad local; los controles de seguridad mantienen sus pruebas y puerta G5. La previsión documental no acredita implementación ni resultados.
\item[CT-04.] S4, Anexo 4-R, protección del canal AS2 mediante firma, cifrado S/MIME y MDN. Este sustento de autenticidad no acredita por sí solo cuarentena, correlación ni rechazo de repeticiones; ARQ/SEG deben completar esa trazabilidad de diseño en S4 y CAL los casos negativos EDI en S9, para E2.
\item[CT-05.] S4, Anexo 4-R, controles 8.7/8.8 y 8.24 y «Ciclo de identidad y custodia»: protección del terminal, cifrado y revocación conectada. SRE/SEG deben concretar en S4 las capacidades de gestión remota disponibles; CAL debe vincular en S9 el ensayo de extravío, sin presumir bloqueo remoto durante la pérdida de enlace.
\item[CT-06.] S4, Anexo 4-R, control 8.20/8.22: borde y superficie B2B separada. T-13, 7.1, define los umbrales de carga, no un ensayo específico de denegación de servicio. SEG/SRE deben completar límites y protección en S4 y CAL vincular en S9 la ráfaga autorizada y las mediciones de la ficha.
\item[CT-07.] S4, Anexo 4-R, controles 8.2, 8.13/8.14 y 8.15/8.16: segregación, respaldo y auditoría; T-13, 7.5, AL-DR-01, aporta recuperación del WMS. No equivale a probar abuso administrativo. SEG/SRE deben completar el escenario de contención/restauración en S4 y CAL su verificación específica en S9. Para CT-04--CT-07 no se acredita aún correspondencia individual con casos de S9; SEG y CAL deben registrarla, sin cerrar los controles con estas referencias generales.
\item[CT-08.] S13, 13.1.6, I-01A, y 13.1.7, «Control de adopción»: validación del catálogo, certificación y piloto antes de ampliar rutas. El indicador y el piloto son evidencia prevista, no aciertos ya observados.
\item[CT-09 y CT-10.] S13, 13.2.6, I-02A/I-02B e I-02D, y 13.2.7, «Controles de adopción y seguridad»: CT-09 cubre fidelidad y consumo; CT-10, sanitización previa a exportar. CAL debe conservar la comparación incidente--reproducción y el consumo medido; SEG/CAL, la prueba negativa de datos prohibidos en cada liberación y cambio de esquema. Los indicadores no sustituyen esos registros.
\item[CT-11 y CT-12.] S13, 13.3.6, I-03A/I-03B, y 13.3.7, «Controles de validación por familia y trazabilidad térmica»: CT-11 valida por familia sin relajar RNG-04 ni extender la fecha impresa; CT-12 ensaya movimiento sin evento, sensor equivocado y sincronización tardía, con historia incompleta visible y sin estimación. DAT/CAL deben aportar los resultados antes de habilitar cada familia.
\item[CT-13.] S13, 13.5.6, I-05B/I-05D, y 13.5.7, «Control de comprensión y privacidad»: comprensión del prototipo, revisión legal y umbral/supresión en cada emisión comparativa. IMP/SEG/CAL deben registrar resultados antes de habilitar los bloques respectivos; la prueba de comprensión no acredita por sí sola privacidad.
\item[CT-14.] S4, «Primer cuello de botella bajo la carga de septiembre», tabla «Detección del primer cuello de botella»: ERP por ACL, base de Talca, cola y M2. T-13, 7.1, especifica cargas, ventana sostenida, crecimiento/quiebre, p95 e informe por componente; 7.5, AL-PERF-01, complementa la carga con aislamiento por perfil. SRE/ARQ/CAL deben contrastar la hipótesis con mediciones, no con el cálculo de dimensionamiento.
\item[CT-15 (stock).] S4, Anexo 4-V, AL-STOCK-01, y Anexo 4-M, AL-DR-01, con procedimientos de conmutación/retorno de 4.3.2: retención local durable, exclusión del escritor previo y conciliación. T-13 de S9, 7.5, asigna 6.1.1 a preventa E1, 6.1.2 a preventa/portal/EDI E2 y 6.4.1/6.4.2 a recuperación. T-17, 3.1, conserva el expediente y el resultado por etapa. Un simulador no acredita producto E2. CT-16 histórico de stock equivale a CT-15 vigente; la identidad local antes separada se integra en CT-03, sin reutilizar CT-15 para identidad.
\end{description}

CAL debe incorporar al expediente aplicable la versión ensayada, entradas, resultados, defectos y acta; JP mantiene pendientes las dependencias de diseño/prueba hasta su resolución en las ramas responsables. Esta trazabilidad documental no acredita implementación, ejecución ni aceptación técnica.

\subsection*{Condiciones de identidad local de CT-03}
El manifiesto de bodega dura 26 h y el ensayo corta el enlace una hora después
de emitirlo: quedan 25 h verificables. La caché descriptiva de 24 h no emite
credenciales ni amplía el manifiesto. En los relevos a las 8 y 16 h del corte se
validan firma, emisor, sitio, dispositivo enrolado, rol, PIN y turno vigente;
sin enlace no se habilitan cuentas no preinscritas ni privilegios nuevos.
Manifiesto o turno vencido, permiso inválido, firma alterada o PIN incorrecto
se deniegan y auditan. Reinicio, actividad o retroceso del reloj no extienden
vigencia; sin continuidad demostrable de la hora confiable se deniega sesión nueva.

En reparto, la autorización se emite dentro de la hora previa a salir, cubre
la ruta de 14,5 h y expira al fin programado más 30 min, con máximo de 16 h
desde su emisión. Al vencer, se deniegan nuevas operaciones protegidas y se
conserva la cola. La reautenticación conectada permite conciliar una sola vez
los hechos anteriores al vencimiento; otra ruta exige autorización nueva.

Una baja conocida en el sitio deniega acceso local inmediato. Una baja conocida
sólo por la autoridad central durante el aislamiento deja exposición hasta la
expiración: no se promete revocación remota instantánea. Al reconectar se bloquea
la sesión y se ponen en cuarentena los eventos posteriores a la baja para
conciliación, sin borrarlos. El expediente de AL-ID-01 conserva emisión,
expiración, huella del manifiesto, política, versión del verificador, tiempos,
decisiones y UUID; excluye PIN y secretos. CAL y SEG verifican resultados y
la aprobación competente habilita sólo el alcance probado.
